% from ega4-2-fr.tex
\oldpage[IV-2]{227}

\EGAIVtwoTerm{$\mathcal O_X$-Module géométriquement ponctuellement intègre sur un corps}{4.6.22}
\EGAIVtwoTerm{$\mathcal O_X$-Module géométriquement réduit (ou séparable) sur un corps}{4.6.17}
\EGAIVtwoTerm{$\mathcal O_X$-Module géométriquement réduit (ou séparable) sur un corps en un point}{4.6.22}
\EGAIVtwoTerm{$\mathcal O_X$-Module équidimensionnel en un point}{5.1.12}
\EGAIVtwoTerm{$\mathcal O_X$-Module intègre, intègre en un point}{3.2.4}
\EGAIVtwoTerm{$\mathcal O_X$-Module irredondant}{3.2.4}
\EGAIVtwoTerm{$\mathcal O_X$-Module normalement plat le long d'un sous-préschéma}{6.10.1}
\EGAIVtwoTerm{$\mathcal O_X$-Module possédant la propriété $(S_k)$,
possédant la propriété $(S_k)$ en un point}{5.7.2}
\EGAIVtwoTerm{$\mathcal O_X$-Module réduit, réduit en un point}{3.2.2}
\EGAIVtwoTerm{$\mathcal O_X$-Module $Z$-clos, $Z$-pur}{5.9.9}
\EGAIVtwoTerm{$\mathcal O_X$-Module $Z$-clos, $Z$-pur en un point}{5.9.9}
\EGAIVtwoTerm{Morphisme birationnel de préschémas réduits}{6.15.4}
\EGAIVtwoTerm{Morphisme connexe}{4.5.5}
\EGAIVtwoTerm{Morphisme de Cohen-Macaulay en un point, morphisme de Cohen-Macaulay}{6.8.1}
\EGAIVtwoTerm{Morphisme de coprofondeur $\leq n$ en un point, morphisme de coprofondeur $\leq n$}{6.8.1}
\EGAIVtwoTerm{Morphisme défini sur un corps}{4.8.4}
\EGAIVtwoTerm{Morphisme irréductible}{4.5.5}
\EGAIVtwoTerm{Morphisme lisse en un point, morphisme lisse}{6.8.1}
\EGAIVtwoTerm{Morphisme normal en un point, morphisme normal}{6.8.1}
\EGAIVtwoTerm{$P$-morphisme}{7.3.1}
\EGAIVtwoTerm{Morphisme possédant la propriété $(R_k)$ en un point,
morphisme possédant la propriété $(R_k)$}{6.8.1}
\EGAIVtwoTerm{Morphisme possédant la propriété $(S_k)$ en un point,
morphisme possédant la propriété $(S_k)$}{6.8.1}
\EGAIVtwoTerm{Morphisme quasi-fidèlement plat, morphisme quasi-plat}{2.3.3}
\EGAIVtwoTerm{Morphisme radiciel en un point}{6.15.3}
\EGAIVtwoTerm{Morphisme réduit en un point, morphisme réduit}{6.8.1}
\EGAIVtwoTerm{Morphisme régulier en un point, morphisme régulier}{6.8.1}
\EGAIVtwoTerm{Morphisme résolvant}{7.9.1}
\EGAIVtwoTerm{Morphisme universellement ouvert, universellement bicontinu}{2.4.2}
\EGAIVtwoTerm{Multiplicité radicielle, multiplicité séparable, multiplicité
totale d'une extension d'un corps}{4.7.4}
\EGAIVtwoTerm{Multiplicité radicielle, multiplicité séparable, multiplicité
totale d'un préschéma sur un corps}{4.7.4}
\EGAIVtwoTerm{Multiplicité radicielle d'un point pour un $\mathcal O_X$-Module
($X$ préschéma sur un corps)}{4.7.5}
\EGAIVtwoTerm{Multiplicité radicielle d'un cycle premier maximal associé à un
$\mathcal O_X$-Module ($X$ préschéma sur un corps)}{4.7.8}
\EGAIVtwoTerm{Multiplicité totale d'un point (d'un cycle premier maximal
associé) pour un $\mathcal O_X$-Module ($X$ préschéma sur un corps)}{4.7.12}

\medskip
\EGAIVtwoTerm{Nombre géométrique de composantes connexes, de composantes
irréductibles, pour un préschéma sur un corps}{4.5.2}

\medskip
\EGAIVtwoTerm{Point associé à un $\mathcal O_X$-Module, à un préschéma}{3.1.1}
\EGAIVtwoTerm{Point de Cohen-Macaulay d'un préschéma}{5.7.1}
\EGAIVtwoTerm{Point géométriquement unibranche, point unibranche d'un préschéma}{6.15.1}
\EGAIVtwoTerm{Préschéma connexe en codimension $d-1$}{5.10.8}
\EGAIVtwoTerm{Préschéma de Cohen-Macaulay}{5.7.1}
\EGAIVtwoTerm{Préschéma excellent}{7.8.5}
\EGAIVtwoTerm{Préschéma (sur un corps) géométriquement connexe}{4.5.2}
\EGAIVtwoTerm{Préschéma (sur un corps) géométriquement intègre}{4.6.2}
\EGAIVtwoTerm{Préschéma (sur un corps) géométriquement ponctuellement intègre}{4.6.9}
\EGAIVtwoTerm{Préschéma (sur un corps) géométriquement irréductible}{4.5.2}
\EGAIVtwoTerm{Préschéma (sur un corps) géométriquement normal en un point,
géométriquement normal}{6.7.6}
\EGAIVtwoTerm{Préschéma (sur un corps) géométriquement réduit}{4.6.2}
\EGAIVtwoTerm{Préschéma (sur un corps) géométriquement réduit en un point}{4.6.9}
\EGAIVtwoTerm{Préschéma (sur un corps) géométriquement régulier en un point,
géométriquement régulier}{6.7.6}

\clearpage
